\documentclass{article}
\usepackage{amsmath}
\usepackage{amssymb}
\usepackage{graphicx} % Required for inserting images

\title{Finite Difference Scheme for the Ornstein–Uhlenbeck process}
\author{Yiyang Hu}
\date{July 2026}

\begin{document}

\maketitle


\noindent Consider the following Fokker--Planck equation defined on
$[x_L,x_R]$, where $x_L=-5$ and $x_R=5$:

$$
\partial_t\rho
=
\kappa\partial_x(x\rho)
+
D\partial_{xx}\rho,
\qquad
x\in[x_L,x_R],
\qquad
t\geq 0.
$$

The probability flux is defined by

$$
J(x,t)
=
-\kappa x\rho(x,t)
-
D\partial_x\rho(x,t).
$$

We impose the zero-flux boundary conditions

$$
J(x_L,t)=J(x_R,t)=0,
\qquad
t\geq 0.
$$

Let $T>0$ be the maximum time. For $M,N\in\mathbb{N}$, define

$$
h=\frac{x_R-x_L}{M}
=
\frac{10}{M},
\qquad
\delta=\frac{T}{N}.
$$

The spatial and temporal grid points are given by

$$
x_j=x_L+jh,
\qquad
j=0,1,\ldots,M,
$$

and

$$
t^{(n)}=n\delta,
\qquad
n=0,1,\ldots,N.
$$

We write

$$
\rho_j^n
:=
\rho(x_j,t^{(n)}).
$$

At the interior grid points, the first and second spatial derivatives
are approximated using central differences:

$$
\partial_x(x\rho)(x_j,t^{(n)})
\approx
\frac{x_{j+1}\rho_{j+1}^n-x_{j-1}\rho_{j-1}^n}{2h},
$$

and

$$
\partial_{xx}\rho(x_j,t^{(n)})
\approx
\frac{\rho_{j+1}^n-2\rho_j^n+\rho_{j-1}^n}{h^2}.
$$

Using the forward Euler method for the time derivative, we obtain

$$
\frac{\rho_j^{n+1}-\rho_j^n}{\delta}
=
\kappa
\frac{x_{j+1}\rho_{j+1}^n-x_{j-1}\rho_{j-1}^n}{2h}
+
D
\frac{\rho_{j+1}^n-2\rho_j^n+\rho_{j-1}^n}{h^2}.
$$

Therefore, for $j=1,\ldots,M-1$,

$$
\rho_j^{n+1}
=
\rho_j^n
+
\frac{\kappa\delta}{2h}
\left(
x_{j+1}\rho_{j+1}^n
-
x_{j-1}\rho_{j-1}^n
\right)
+
\frac{D\delta}{h^2}
\left(
\rho_{j+1}^n
-
2\rho_j^n
+
\rho_{j-1}^n
\right).
$$

\section*{Elimination of the ghost points}

To encode the boundary conditions, introduce the ghost points
$x_{-1}=x_0-h$ and $x_{M+1}=x_M+h$.

At the left boundary, the zero-flux condition gives

$$
0
=
-\kappa x_0\rho_0^n
-
D\frac{\rho_1^n-\rho_{-1}^n}{2h}.
$$

Hence,

$$
D\frac{\rho_1^n-\rho_{-1}^n}{2h}
=
-\kappa x_0\rho_0^n,
$$

and therefore

$$
\rho_{-1}^n
=
\rho_1^n
+
\frac{2\kappa h}{D}x_0\rho_0^n.
$$

Similarly, at the right boundary,

$$
0
=
-\kappa x_M\rho_M^n
-
D\frac{\rho_{M+1}^n-\rho_{M-1}^n}{2h},
$$

so that

$$
\rho_{M+1}^n
=
\rho_{M-1}^n
-
\frac{2\kappa h}{D}x_M\rho_M^n.
$$

For notational convenience, define

$$
\alpha
:=
\frac{\kappa\delta}{2h},
\qquad
\beta
:=
\frac{D\delta}{h^2}.
$$

We also define the boundary coefficients

$$
\begin{aligned}
a_L
&:=
-\frac{\kappa^2\delta}{D}x_{-1}x_0,
&
a_R
&:=
-\frac{\kappa^2\delta}{D}x_{M+1}x_M,
\\[6pt]
b_L
&:=
\frac{2\kappa h}{D}x_0-2,
&
b_R
&:=
-\frac{2\kappa h}{D}x_M-2.
\end{aligned}
$$

Let the discrete density vector be

$$
\boldsymbol{\rho}^{\,n}
:=
\begin{pmatrix}
\rho_0^n &
\rho_1^n &
\cdots &
\rho_M^n
\end{pmatrix}^{\mathsf T}.
$$

After eliminating the ghost values, the fully discrete scheme can be
written in matrix form as

$$
\boldsymbol{\rho}^{\,n+1}
=
\left(
I_{M+1}+A+B
\right)
\boldsymbol{\rho}^{\,n},
$$

where $I_{M+1}$ is the $(M+1)\times(M+1)$ identity matrix,
$A$ represents the drift contribution, and $B$ represents the
diffusion contribution.

\begingroup

\renewcommand{\arraystretch}{1.3}
\setlength{\arraycolsep}{7pt}

The drift matrix is

$$
A
=
\begin{pmatrix}
a_L
&
\kappa\delta
&
0
&
\cdots
&
0
\\
-\alpha x_0
&
0
&
\alpha x_2
&
\ddots
&
\vdots
\\
0
&
-\alpha x_1
&
0
&
\ddots
&
0
\\
\vdots
&
\ddots
&
\ddots
&
\ddots
&
\alpha x_M
\\
0
&
\cdots
&
0
&
\kappa\delta
&
a_R
\end{pmatrix}.
$$

The diffusion matrix is

$$
B
=
\beta
\begin{pmatrix}
b_L
&
2
&
0
&
\cdots
&
0
\\
1
&
-2
&
1
&
\ddots
&
\vdots
\\
0
&
1
&
-2
&
\ddots
&
0
\\
\vdots
&
\ddots
&
\ddots
&
\ddots
&
1
\\
0
&
\cdots
&
0
&
2
&
b_R
\end{pmatrix}.
$$

\endgroup

Thus, the numerical solution is updated at each time step according to

$$
\boxed{
\boldsymbol{\rho}^{\,n+1}
=
\left(
I_{M+1}+A+B
\right)
\boldsymbol{\rho}^{\,n}
}.
$$

\end{document}

